\documentclass[10pt,twocolumn]{article}
\usepackage[margin=0.75in]{geometry}
\usepackage{amsmath,amssymb,amsthm}
\usepackage{booktabs}
\usepackage{graphicx}
\usepackage[hidelinks]{hyperref}
\usepackage{enumitem}
\usepackage{microtype}
\usepackage{caption}
\usepackage{dblfloatfix}
\captionsetup{font=small,labelfont=bf}

\newtheorem{definition}{Definition}
\newtheorem{remark}{Remark}

\setlength{\parskip}{0.3em}
\setlength{\columnsep}{0.3in}

\title{\textbf{Execution Cost Lab}\\[4pt]
\large TWAP vs Forecast-Profile VWAP\\on Intraday Equity Bars\\
under Participation Constraints}

\author{Alice Lemaire\\
CentraleSup\'elec \& EPFL (Double Degree)\\
\texttt{alice.lemaire@student.centralesupelec.fr}}

\date{April 2026}

\begin{document}
\maketitle

%% ═══════════════════════════════════════════════
\begin{abstract}
When an institutional investor executes a large buy order, the order must be split into smaller slices spread over time to avoid moving the price. The two standard scheduling policies are TWAP (uniform slicing) and VWAP (volume-weighted slicing). This paper builds a simulation framework to quantify whether a historical intraday volume profile improves execution cost relative to a na\"ive TWAP.

I introduce a three-way decomposition---TWAP, forecast-profile VWAP (tradeable), and realized-volume VWAP (oracle, not tradeable)---that separates the value of volume information from the cost of forecast error. Over 1{,}000 simulated 60-minute windows on five US equities, the forecast VWAP wins 50.9\% of windows on implementation shortfall. A regime analysis reveals that this advantage is concentrated in high-volume environments and reverses in low-volume, high-volatility conditions. The full pipeline is implemented as a companion Jupyter notebook.
\end{abstract}

%% ═══════════════════════════════════════════════
\section{Introduction}

Suppose a portfolio manager decides to buy 50{,}000 shares of NVDA. If the full order is sent at once, it would consume all available liquidity at the current best ask price, then sweep through deeper levels of the order book, pushing the price higher with each fill. This price degradation---called \emph{market impact}~\cite{bouchaud2009}---is not a theoretical concern: Berkowitz et al.~\cite{berkowitz1988} estimated that execution costs on the NYSE can reach 1--2\% of trade value, far exceeding explicit commissions.

To mitigate impact, the order is split into smaller child orders spread over a time horizon~\cite{cartea2015}. The question becomes: \emph{how should these slices be distributed?} This is the core problem of algorithmic execution.

I compare three scheduling strategies: TWAP (uniform), forecast-profile VWAP (based on a historical volume curve), and oracle VWAP (based on actual realized volume, not implementable in practice). The paper makes three contributions:
\begin{enumerate}[nosep]
\item A clean separation between the tradeable forecast and the non-tradeable oracle, with the gap measuring the \textbf{cost of forecast error}---a diagnostic that indicates whether effort should go into better forecasting or into other aspects of execution.
\item A \textbf{regime analysis} showing that the VWAP advantage is not uniform across market conditions: it is strongest in high-volume environments and can reverse when volume is low and volatility is high.
\item A \textbf{fully reproducible pipeline} (Jupyter notebook, fixed seed) where every number and figure in this paper can be regenerated.
\end{enumerate}

All simulations are buy-side: positive implementation shortfall means I paid more than the arrival price (unfavourable).


%% ═══════════════════════════════════════════════
\section{Background}

This section explains two concepts that are essential to understanding the rest of the paper.

\subsection{The participation cap}

When an algorithm executes a buy order, it must not consume too large a fraction of the available volume on any given bar. If the algorithm buys 30\% of a bar's total volume, this is highly visible to other market participants and creates significant price impact. In practice, execution algorithms cap their participation at 5--20\% of each bar's volume.

In this paper, I use a cap of $\rho = 5\%$: at any bar~$t$ with market volume $V_t$, the algorithm can execute at most $0.05 \times V_t$ shares. This is a conservative, realistic constraint. If the algorithm's desired quantity exceeds this cap, part of the order is left unexecuted on that bar.

\subsection{The U-shaped intraday volume pattern}

Trading volume is not uniformly distributed throughout the day. On most equity exchanges, volume follows a distinctive U-shape~\cite{admati1988, biais1995}: high volume at the open (9:30--10:00, driven by overnight information and institutional order flow), low volume at midday (12:00--14:00), and high volume again at the close (15:00--16:00, driven by portfolio rebalancing and index tracking).

This pattern is the central motivation for VWAP scheduling. If the algorithm knows the volume shape in advance, it can concentrate execution during high-volume periods (where it consumes a smaller fraction of available liquidity) and reduce execution during quiet periods (where even a small order creates a large footprint).



%% ═══════════════════════════════════════════════
\section{Data}

\subsection{Universe}
I use 5-minute OHLCV bars for five instruments: AAPL, MSFT, AMZN, NVDA (mega-cap tech stocks) and SPY (S\&P~500 ETF), downloaded via \texttt{yfinance} over 60 trading days. This universe offers a range of liquidity profiles: SPY has extremely deep, relatively flat intraday volume; NVDA has higher volatility and a more pronounced U-shaped volume curve.

\subsection{Cleaning}
Raw timestamps arrive in UTC with mixed timezone offsets across DST transitions (EST \texttt{-05:00} in winter, EDT \texttt{-04:00} in summer). The pipeline: (1)~converts to US/Eastern time; (2)~deduplicates; (3)~keeps weekdays only; (4)~filters to regular trading hours 09:30--16:00~ET; (5)~removes zero-volume bars. Step~(1) is critical because the volume profile is defined relative to local trading hours. Step~(5) prevents division-by-zero in VWAP computations.

After cleaning, the dataset contains approximately 4{,}500 bars per ticker, spanning $\sim$60 trading sessions over roughly three calendar months. Each session has 78 bars (one per 5-minute interval from 09:30 to 16:00).


%% ═══════════════════════════════════════════════
\section{Model}

\subsection{Setup}

A buy order of $Q$ shares is executed over $N\!=\!12$ consecutive 5-minute bars (one hour). The order size is set as 3\% of the total volume traded during the window:
\begin{equation}
Q = 0.03 \times \sum\nolimits_{t=1}^{N} V_t\,.
\end{equation}
This is a moderate order size. Since the participation cap allows up to $5\%$ of total volume, full execution is feasible in aggregate---but uneven volume distribution across bars can prevent full fill in practice.

\begin{remark}[Design choices]
$\alpha = 0.03$ represents a moderate institutional order (3\% of the window's total volume). $\rho = 0.05$ is a conservative participation cap, typical of institutional execution algorithms. $N = 12$ (one hour) is a standard execution horizon for medium-sized orders. The no-catch-up rule is deliberate: reallocating unfilled shares to later bars would confound the scheduling effect with the reallocation heuristic, making it impossible to attribute performance differences to the schedule alone.
\end{remark}

\subsection{Shared execution rule}

All three strategies compute a \emph{desired quantity} $d_t$ at each bar, then execute the actual quantity:
\begin{equation}
q_t = \min(d_t,\;\rho\, V_t,\;Q_{\text{rem}})\,,
\label{eq:exec}
\end{equation}
where $Q_{\text{rem}} = Q - \sum_{s<t} q_s$ is the unfilled remainder. The three terms enforce, respectively: (i)~the strategy's scheduling rule, (ii)~the participation cap (never exceed 5\% of the bar's volume), and (iii)~the no-overfill constraint (stop once the full order is filled). Execution is priced at each bar's close price~$P_t$.

Importantly, there is \emph{no catch-up}: if the cap prevents filling bar~$t$'s target, the shortfall is not reallocated to subsequent bars. This isolates the schedule's quality from reallocation heuristics.

\subsection{Strategy 1: TWAP}

The desired quantity is uniform:
\begin{equation}
d_t^{\,\text{TWAP}} = Q / N\,.
\end{equation}
TWAP is simple and requires no market information. Its weakness: it attempts to trade the same quantity during a quiet midday bar (where volume may be 50{,}000 shares) as during a busy opening bar (where volume may be 500{,}000 shares). On the quiet bar, the participation cap binds and the order is partially unfilled.

\subsection{Strategy 2: Realized-volume VWAP (oracle)}

The desired quantity is proportional to the bar's actual volume:
\begin{equation}
d_t^{\,\text{oracle}} = Q \cdot V_t \,/\, {\textstyle\sum_s V_s}\,.
\end{equation}
This means: trade a lot when the market is active, trade little when it is quiet. The participation rate $q_t / V_t$ is constant across bars (before the cap binds), which is the theoretically ideal behaviour.

However, this schedule requires knowing $V_t$ \emph{before} bar~$t$ completes---information that is not available in real time. It is therefore \textbf{not tradeable} and serves purely as an oracle: the best any volume-proportional schedule could achieve with perfect foresight.

\subsection{Strategy 3: Forecast-profile VWAP}

To build a tradeable VWAP, I estimate the intraday volume shape from historical data.

\begin{definition}[Intraday volume profile]
For each past trading session~$d$, compute the fraction of daily volume at each bar index~$k$: $\varphi_{d,k} = V_{d,k}/\sum_s V_{d,s}$. Average across sessions: $\bar\varphi_k = D^{-1}\sum_d \varphi_{d,k}$. This curve captures the U-shape.
\end{definition}

For a window starting at bar index $k_0$, the desired quantity is:
\begin{equation}
d_t^{\,\text{fcst}} = Q \cdot \frac{\bar\varphi_{k_0+t-1}}{\sum_{s}\bar\varphi_{k_0+s-1}}\,.
\end{equation}

In words: bars whose historical profile predicts higher volume receive a proportionally larger slice of the order. The profile is built from sessions \textbf{strictly before} the window's date, ensuring the forecast is genuinely ex-ante.

\subsection{Metrics}

\textbf{Implementation shortfall} (IS), introduced by Perold~\cite{perold1988}, is the standard metric for execution quality. It measures the difference between the average price actually paid and the price at which the decision was made (the \emph{arrival price}, here defined as the first bar's close):
\begin{equation}
\text{IS} = \bar{P}_{\text{exec}} - P_0, \quad \bar{P}_{\text{exec}} = \frac{\sum_t q_t P_t}{\sum_t q_t}\,.
\label{eq:IS}
\end{equation}
For a buy order, IS~$> 0$ means the market moved up during execution---I paid more than at decision time (unfavourable). IS~$< 0$ means the market moved down during execution (favourable). IS captures all implicit costs---impact, timing, drift---in a single number expressed in USD/share. It is computed on executed shares only.

\textbf{Fill ratio}: $\text{Fill} = \sum_t q_t / Q$. Values below~1 indicate that the participation cap prevented full execution. Fill is tracked separately from IS to avoid conflating schedule quality with underfill: a strategy that fills only 80\% of the order at a good price should not be confused with one that fills 100\% at a bad price.


%% ═══════════════════════════════════════════════
\section{Experimental Design}

\subsection{Pipeline structure}

The codebase consists of six Python modules executed sequentially: \texttt{download\_data} (fetch OHLCV from Yahoo Finance), \texttt{clean\_data} (timezone conversion and filtering), \texttt{execution} (TWAP/VWAP simulators and metric computation), \texttt{run\_experiment} (window sampling and strategy comparison), \texttt{make\_plots} (summary tables and figures), and \texttt{run\_one\_window} (single-window visual demo). A shell script \texttt{run\_all.sh} orchestrates the full pipeline end-to-end.

\subsection{Window sampling}

For each ticker, I enumerate all valid contiguous 12-bar windows within each trading session and sample 200 at random (seed~=~42), yielding 1{,}000 windows across five tickers. All bars in a window come from a single session---mixing sessions would break the arrival price reference and distort the volume profile.

Each window is characterised by two conditioning variables used in the regime analysis: \textbf{realized volatility} $\sigma = \sqrt{\sum r_t^2}$ (from 5-minute log-returns) and \textbf{volume CV} $= \sigma_V/\mu_V$ (how unevenly volume is distributed across the 12 bars). High CV means some bars are much more active than others---exactly the condition where VWAP should outperform TWAP.


%% ═══════════════════════════════════════════════
\section{Results}

\subsection{Single-window illustration}

Figure~\ref{fig:demo} shows one 60-minute window on AAPL. In the top panel, TWAP attempts a constant slice at every bar; on quiet bars (centre of the window), the participation cap (dotted line) binds and TWAP cannot fill its target. The forecast VWAP, by contrast, assigns less to quiet bars and more to active ones, avoiding cap violations. The oracle matches the volume shape perfectly. The middle panel shows cumulative fill; the bottom panel shows the price path.

\begin{figure}[t]
\centering
\includegraphics[width=\linewidth]{single_window_demo_AAPL.png}
\caption{Single-window demo (AAPL). Top: per-bar execution vs market volume. Middle: cumulative fill. Bottom: price path.}
\label{fig:demo}
\end{figure}

\subsection{Aggregate results}

Table~\ref{tab:ticker} reports per-ticker averages over the 1{,}000 windows.

\begin{table}[t]
\centering
\caption{Per-ticker summary (IS in USD/share).}
\label{tab:ticker}
\small
\begin{tabular}{@{}lrrrrr@{}}
\toprule
 & TWAP & Fcst & TWAP & Fcst & Win \\
 & IS & IS & fill & fill & \% \\
\midrule
AAPL & $-$.088 & $-$.082 & .958 & .956 & 46.5 \\
AMZN & $+$.053 & $+$.052 & .968 & .965 & 50.5 \\
MSFT & $-$.090 & $-$.082 & .943 & .943 & 52.5 \\
NVDA & $-$.046 & $-$.052 & .944 & .930 & 53.0 \\
SPY  & $+$.036 & $+$.034 & .992 & .994 & 52.0 \\
\bottomrule
\end{tabular}
\end{table}

Win rates range from 46.5\% (AAPL) to 53.0\% (NVDA). The advantage is real but modest.

\textit{A note on negative IS values.} Several tickers show negative mean IS (e.g., AAPL: $-$0.088, MSFT: $-$0.090). This does \emph{not} mean the strategies are generating profit---it means the market happened to move down on average during the sampled windows while the buy order was executing, a favourable coincidence driven by the sampled price paths. The meaningful comparison is always the \emph{difference} $\Delta$IS between strategies, which cancels out the common price drift.

Three observations:

\textit{NVDA has the highest win rate.} NVDA has the most pronounced U-shaped volume pattern among the five tickers. When the shape is strong, the historical average captures it well and the forecast VWAP benefits.

\textit{AAPL is below 50\%.} This suggests that Apple's intraday volume pattern during the sample period was either flat (little to exploit) or unstable (the average profile is a poor forecast of any given day's shape).

\textit{Fill ratios are slightly lower for the forecast VWAP.} By concentrating execution in specific bars, the forecast VWAP is more exposed to cap violations when the forecast is wrong. The exception is SPY, where liquidity is so deep that the cap almost never binds.

\subsection{The IS difference distribution}

Figure~\ref{fig:hist} shows $\Delta\text{IS} = \text{IS}_{\text{fcst}} - \text{IS}_{\text{TWAP}}$ per window. The distribution is centred near zero with mean $\approx -0.003$ USD/share. The wide spread ($\pm\,0.5$) reveals an important insight: on any single window, whether the market moved up or down during the hour matters far more than the scheduling choice. The scheduling effect is a second-order improvement on top of first-order price randomness.

\begin{figure}[t]
\centering
\includegraphics[width=\linewidth]{hist_delta_is_forecast_vwap_vs_twap.png}
\caption{$\Delta$IS distribution. Green = forecast wins.}
\label{fig:hist}
\end{figure}

\subsection{Regime analysis: when does VWAP help?}

This is the central empirical finding. I split the 1{,}000 windows into four quadrants by the median of window volume ($V_{60}$) and realized volatility ($\sigma$). Table~\ref{tab:regime} reports the results.

\begin{table}[t]
\centering
\caption{Regime analysis.}
\label{tab:regime}
\small
\begin{tabular}{@{}llrcr@{}}
\toprule
Volume & Volatility & $n$ & $\overline{\Delta\text{IS}}$ & Win\% \\
\midrule
High & High & 282 & $-$.0003 & 54.3 \\
High & Low  & 218 & $-$.0030 & 50.9 \\
Low  & High & 218 & $+$.0070 & 46.3 \\
Low  & Low  & 282 & $+$.0013 & 51.1 \\
\bottomrule
\end{tabular}
\end{table}

The pattern is clear. In \textbf{high-volume} environments (top two rows), the forecast VWAP outperforms TWAP consistently (54.3\% and 50.9\% win rates). In the \textbf{low-volume, high-volatility} quadrant, TWAP is better (only 46.3\% win rate for the forecast, $\Delta\text{IS} = +0.007$).

The interpretation is intuitive. When volume is high, the intraday pattern is strong---there is a clear difference between busy and quiet bars---and the historical profile captures this well. The forecast VWAP exploits the pattern by loading execution into liquid bars. When volume is low and volatility is high, the pattern is noisy: volume on any given bar is hard to predict, and concentrating execution based on a poor forecast is worse than spreading it uniformly.

Figure~\ref{fig:regime} visualises this as a heatmap.

\begin{figure}[t]
\centering
\includegraphics[width=\linewidth]{regime_heatmap.png}
\caption{Regime heatmap. Left: mean $\Delta$IS. Right: win rate. High-volume rows favour the forecast.}
\label{fig:regime}
\end{figure}

\subsection{Oracle decomposition}

The three-way comparison (Figure~\ref{fig:meanis}) enables a diagnostic decomposition that is the paper's key conceptual contribution:

\begin{itemize}[nosep]
\item \textbf{Oracle vs TWAP}: the maximum possible benefit of knowing the volume shape perfectly. This is an upper bound---no real strategy can exceed it.
\item \textbf{Forecast vs TWAP}: the actual benefit delivered by the historical average profile.
\item \textbf{Forecast vs Oracle}: the cost of forecast error. If this gap is large, a better forecast (exponential weighting, day-of-week conditioning) would help. If it is small, the forecast is already close to the best possible.
\end{itemize}

For NVDA, the oracle substantially beats TWAP and the forecast captures most of the gap: the volume pattern is strong \emph{and} stable over time, so the historical average is a good predictor. For AAPL, the oracle advantage is smaller and the forecast misses most of it: the pattern is either weak or unstable.

This decomposition is actionable. If the oracle-to-TWAP gap is small for a given ticker, scheduling is not the bottleneck---the priority should be reducing other execution costs (spread, impact). If the gap is large but the forecast captures little of it, investing in better volume forecasting would be the highest-return improvement.

\begin{figure}[t]
\centering
\includegraphics[width=\linewidth]{mean_is_by_ticker.png}
\caption{Mean IS by strategy. The forecast-to-oracle gap is the cost of forecast error.}
\label{fig:meanis}
\end{figure}

\subsection{Supporting figures}

Figure~\ref{fig:boxplot} shows IS distributions. The within-strategy spread (interquartile range) is much larger than the between-strategy gap, confirming that market conditions dominate scheduling effects.

Figure~\ref{fig:fill} shows fill ratios. The forecast VWAP fills slightly less than TWAP on all tickers except SPY, reflecting the cap-violation trade-off discussed above.

Figure~\ref{fig:winrate} shows win rates by ticker, with the 50\% baseline. No ticker is dramatically above or below, confirming that the effect is systematic but small.

\begin{figure}[t]
\centering
\includegraphics[width=\linewidth]{is_boxplot_by_ticker.png}
\caption{IS distributions. The IQR dwarfs the strategy gap.}
\label{fig:boxplot}
\end{figure}

\begin{figure}[t]
\centering
\includegraphics[width=\linewidth]{fill_ratio_by_ticker.png}
\caption{Fill ratios by ticker.}
\label{fig:fill}
\end{figure}

\begin{figure}[t]
\centering
\includegraphics[width=\linewidth]{forecast_vwap_better_rate_by_ticker_is.png}
\caption{Win rate by ticker (dashed = 50\% baseline).}
\label{fig:winrate}
\end{figure}


%% ═══════════════════════════════════════════════
\section{Discussion}

\subsection{Why the effect is small}

A 50.9\% win rate and a 0.003 USD/share mean advantage may seem disappointing. This is expected. In this simulation, execution is priced at bar close and there is no explicit market impact model. The only way VWAP can outperform is by interacting with the price--volume correlation within each window. This effect is inherently noisy because 60-minute price paths are dominated by random walk behaviour at the 5-minute scale.

An explicit impact model---even a simple one where execution cost is proportional to participation rate---would amplify the VWAP advantage significantly. Under such a model, TWAP's habit of pushing into low-volume bars (where it consumes a large fraction of available liquidity) would be directly penalised, while VWAP's alignment with volume would be rewarded. This is the most impactful extension of the current framework.

\subsection{The fill ratio trade-off}

The forecast VWAP's slightly lower fill ratio (e.g., 93.0\% vs 94.4\% on NVDA) is not a flaw---it is a structural consequence of the strategy. By concentrating execution, VWAP is ``betting'' on the forecast being correct. When it is right (high-volume regimes), the concentration pays off. When it is wrong (low-volume bars that the forecast expected to be busy), the cap binds and the order is underfilled. TWAP, by spreading uniformly, makes no forecast bet and is therefore more robust to forecast error---at the cost of never exploiting the volume shape.

\subsection{Practical implications}

The regime analysis points toward an adaptive strategy: use VWAP during periods when volume is expected to be high (first and last hours of the trading day, earnings weeks) and switch to TWAP when volume is low or uncertain (midday, volatile sessions, illiquid names). An algorithm that monitors real-time volume CV and switches scheduling rule accordingly could capture the VWAP advantage when it exists while avoiding the conditions where it underperforms.


%% ═══════════════════════════════════════════════
\section{Limitations and Extensions}

\textbf{Limitations.}
\begin{enumerate}[nosep]
\item All trades priced at bar close---no intra-bar price dynamics, limit-order simulation, or bid-ask spread.
\item No market impact model; the participation cap is the sole proxy.
\item No catch-up for underfilled bars---a real algorithm would redistribute remaining shares.
\item IS is computed on executed shares only, which can mask underfill.
\item The volume forecast is a plain historical average, with no exponential decay, rolling window, or day-of-week conditioning.
\item No explicit transaction costs (commissions, exchange fees).
\end{enumerate}

\textbf{Extensions.}
\begin{enumerate}[nosep]
\item Add a temporary impact term proportional to participation rate and study how the optimal schedule changes.
\item Replace the historical average with an exponentially-weighted forecast and measure the reduction in the forecast-to-oracle gap.
\item Add a Percentage-of-Volume (POV) strategy that targets a constant participation rate as a third benchmark.
\item Build an adaptive TWAP/VWAP switcher conditioned on real-time volume~CV.
\item Test with sell orders, larger order sizes, and varying participation caps.
\item Apply to futures, FX, or less liquid equities where volume patterns may differ.
\end{enumerate}


%% ═══════════════════════════════════════════════
\section{Conclusion}

I presented a reproducible comparison of three execution schedules---TWAP, forecast-profile VWAP, and oracle VWAP---on intraday equity bars. The forecast VWAP achieves a modest but systematic advantage over TWAP, concentrated in high-volume regimes where the intraday volume pattern is strong and predictable. In low-volume, high-volatility conditions, TWAP is preferable.

The oracle decomposition provides an actionable diagnostic: for each ticker, it quantifies how much value lies in perfect volume information and how much of that value the current forecast captures. In this framework, both quantities are small, suggesting that the most impactful next step is not refining the volume forecast but incorporating explicit market impact into the simulation.

The full pipeline is available as a Jupyter notebook. Every number and figure in this paper can be reproduced.

\begin{thebibliography}{8}
\bibitem{bouchaud2009} J.-P.~Bouchaud, J.~D.~Farmer, F.~Lillo. How markets slowly digest changes in supply and demand. In \textit{Handbook of Financial Markets}, North-Holland, 2009.
\bibitem{cartea2015} \'A.~Cartea, S.~Jaimungal, J.~Penalva. \textit{Algorithmic and High-Frequency Trading}. Cambridge Univ.\ Press, 2015.
\bibitem{perold1988} A.~F.~Perold. The implementation shortfall: paper versus reality. \textit{J.\ Portfolio Management}, 14(3):4--9, 1988.
\bibitem{admati1988} A.~R.~Admati, P.~Pfleiderer. A theory of intraday patterns: volume and price variability. \textit{Rev.\ Financial Studies}, 1(1):3--40, 1988.
\bibitem{biais1995} B.~Biais, P.~Hillion, C.~Spatt. An empirical analysis of the limit order book in the Paris Bourse. \textit{J.\ Finance}, 50(5):1655--1689, 1995.
\bibitem{konishi2002} H.~Konishi. Optimal slice of a VWAP order. \textit{J.\ Financial Markets}, 5(2):197--221, 2002.
\bibitem{berkowitz1988} S.~A.~Berkowitz, D.~E.~Logue, E.~A.~Noser. The total cost of transactions on the NYSE. \textit{J.\ Finance}, 43(1):97--112, 1988.
\bibitem{kissell2003} R.~Kissell, M.~Glantz. \textit{Optimal Trading Strategies}. AMACOM, 2003.
\end{thebibliography}

\end{document}
